\scanpage{40}
Enfin, supposons maintenant que $K$ soit isomorphe $\mathbf R$ [on peut supposer $K=\mathbf R$, car $\mathbf R$ n'a pas d'autom. non trivial] soit $K'$ une clôture algébrique de $K$, $\pi=\mathbf Z/2\mathbf Z$ son groupe de Galois, $X$ lisse \deleted{\unsure{propre}} [projectif] sur $K$. Alors on peut appliquer ce qui précède à $X_{K'}$ sur $K'$. On trouve
\deleted{\[
 H^*(X_{\unsure{K'\mathrm{top}}},\mathbf R)\unsure{\xrightarrow{\sim}}H^*(X)\otimes_K K',
\]}
\[
 [\varphi_{K'}:\coprod_{p,q}H^q(X,\Omega_X^p)\otimes_K K'
 \xrightarrow{\sim}H^*(X_{\deleted{\unsure{K}}})\otimes_K K']
\]
et d'après ce qui a été vu plus haut, cet isomorphisme \uline{commute à} $\pi$, donc provient d'un isomorphisme
\[
 \varphi_{\unsure{i|K}}:\coprod_{p,q}H^q(X,\Omega_X^p)\xrightarrow{\sim}H^*(X).
\]

\margin{On a aussi un $\mathbf R$-isom canonique $(\pi=\mathbf Z/2\mathbf Z)$
\[
 H^*(X_{\mathbf C\mathrm{top}},\mathbf R)\otimes_{\mathbf R}\mathbf C
 \xrightarrow{\sim}H^*(X)\otimes_{\mathbf R}\mathbf C,
\]
où $\pi$ opère : à gauche via son opération sur $X_{\mathbf C\mathrm{top}}$ \uline{et} sur $\mathbf C$ dans $\otimes_{\mathbf R}\mathbf C$
d'où en passant aux invariants sous $\pi$ une décompos. canonique
\[
 H^*(X)\simeq H^*(X)^+\oplus H^*(X)^-,\qquad
 H^*(X)^+\simeq H^*(X_{\mathbf C\mathrm{top}},\mathbf R)^\pi
 =H^*(X_{\mathrm{top}}=X_{\mathbf C\mathrm{top}}/\pi;\mathbf R),
\]
\deleted{$H^{\unsure{*}}(\unsure{X},\mathbf R)$}
\[
 H^*(X)^-=H^*(X_{\mathbf C\mathrm{top}},\unsure{\mathbf R})^-.
\]}

En résumé, chaque fois que $X$ est projectif non singulier sur un corps \deleted{\unsure{top.}} top. loc. cpt complet archimédien, on a une décomposition $\varphi$ en une somme de $H^*(X)$, compatible avec la filtration. \deleted{Si} On pose
\[
 \boxed{H^{p,q}(X)=\varphi(H^q(X,\Omega_X^p))\subset F^pH^q(X)}
\]
Si $K$ n'est pas topologique, alors pour tout choix d'une topologie \unsure{formée|fermée} archimédienne $T$ sur $K$, on obtient sur
\mspage{39}
